\section{Build one working program through separate observations}
Copy \path{drafts/pcm-wav-construction/main.asm} to your working source. The caller
selects bound, parse, filter, encode, load, save, the connected pipeline or the earlier
FIR kernel. Each BEGIN/END region initially returns a safe unfinished result. Replace
that region with your instructions while retaining the supplied labels, saved-register
frames and routes. Use the eight-byte result slots: the name \texttt{count} does not
make every result a byte count.

From the construction directory, run \texttt{python3 prepare.py s01-1 --source working.asm --output observation.asm},
using fresh output names and the recipe named in each step. Preparation replaces only
fixture data and caller controls. It preserves your authored instructions and refuses
to replace an existing output. Continue editing your working source, rather than
restarting from the untouched starter at each observation.

A completed WAV entry returns statuses zero through three. During construction,
status four and zero count mean that execution reached an unfinished region. Earlier
regions may already have written an output prefix or accepted file bytes. Stop at the
named checkpoint before its instruction to inspect that earlier work. Each supplied
after-label has its own \texttt{nop}, so a jump to a different handler does not appear
to have completed the preceding body at the same address. A direct
\texttt{function\_return} observation corresponds to the selected caller's return label;
the pipeline instead stops at \texttt{results\_ready}, before publishing result slots.

The earlier FIR entry has a different result convention: \texttt{v0} is its processed
sample count, not a WAV status. Its missing-body return is minus one with zero clips.
The supplied WAV wrapper compares a positive requested block count with the actual
kernel count and routes an incomplete kernel to development status four. Retain this
guard while bringing forward your earlier routine.

\section{Follow the illustrated private state without changing the public contract}
The parser uses a 160-byte frame, the filter 512 bytes, and encoder, loader and saver
128 bytes each. Saved registers, frame pointer and return address occupy bytes zero
through seventy-nine. The parser's candidate frame count is at eighty. The filter's
input window begins at 112 and its private output at 256. The loader's one-byte probe
is at 112; the saver's private descriptor begins at eighty. None of these temporary
objects is a caller-visible result.

The supplied file descriptor marker in \texttt{s3} starts at minus one. Development
cleanup closes a nonnegative acquired descriptor once. After a completed close, set
the marker back to minus one before entering another operation that can fail or return
unfinished. Nested-call guards propagate failed validation before using its results.
Keep full saved-register values and the original stack pointer through every route.
These register and frame choices support the illustrated guided construction;
equivalent complete solutions may use different private layouts.

Preparing an observation does not create host files or force short reads in an ordinary
IDE run. Recipes declare separate in-memory grading files, permissions, transfer sizes
and descriptor values for the checking system. For a real successful file run, create
\texttt{input.wav} in the launch's working directory with the fixed fixture helper.
Use \texttt{--fixture trace} for the short input and \texttt{--fixture listen} for the
final two-second observation. Forced failures and exact short transfers need the
checking fixture.
